%! Author = sriadityavedantam
%! Date = 3/25/26

\section{Subsequences and the Bolzano–Weierstrass Theorem}

\lesson{9}{Wednesday 3 September 2025 9:10}{Day 9}

\begin{example}[Examples of Sequences]
    \begin{align}
        x_n &= n, \quad \text{or} && 1,2,3,\ldots.\nonumber\\
        x_n &= \dfrac{1}{n},\quad\text{or}&& 1,\dfrac{1}{2},\dfrac{1}{3},\ldots.\label{eq:8.1.1}\\
        x_n &= (-1)^n,\quad\text{or}&& -1,1,-1,1,\ldots.\label{eq:8.1.2}\\
        x_n &= (-1)^n + \dfrac{1}{n}, \quad\text{or}&& 0,\dfrac{3}{2}, -\dfrac{2}{3},\dfrac{5}{4},\ldots.\label{eq:8.1.3}\\
        x_n &= n - \mu(n), \quad\text{or}&&0,1,0,1,2,0,1,2,3,\ldots.\nonumber\\
        \mu(n)&\coloneqq \max\left\{ {m\choose 2} : {m\choose 2}\leq n \right\}.\nonumber\\
        x_n &= \dfrac{1}{n+1 - \mu(n)},\quad\text{or}&&1,\dfrac{1}{2},1,\dfrac{1}{2},\dfrac{1}{3},\ldots.\nonumber\\
        x_n&=\dfrac{n - {m\choose 2}}{m}, \quad {m-1\choose 2}< n\leq {m\choose 2},\ m \geq 2,\nonumber\\
        &\text{ where } {m\choose 2}\coloneqq\begin{cases}
                                                 \dfrac{m(m-1)}{2}, & m\geq 2\\
                                                 0, & m = 1,
        \end{cases}\quad\text{or} && \dfrac{1}{2},\dfrac{1}{3},\dfrac{2}{3},\dfrac{1}{4},\dfrac{1}{2},\dfrac{3}{4},\ldots.\label{eq:8.1.4}
    \end{align}
    Both $\limsup$ and $\liminf$ are limits of convergent subsequences.
\end{example}

\begin{theorem}[Bolzano-Weierstrass]{
    Any bounded sequence of real numbers admits a convergent subsequence.
}
    For each $n\in\n$, define $y_n \coloneqq \sup\{x_m : m\geq n\}$.
    Then $(y_n)$ is monotonically decreasing and bounded below.
    Define
    \[
        \limsup x_n \coloneqq \lim y_n.
    \]
    Define $\liminf x_n$ similarly.
    Suppose $L$ is the limit of some subsequence $(x_{n_k})$.
    Then
    \[
        \liminf x_n \leq L \leq \limsup x_n.
    \]
    Let $L^\upstar \coloneqq \limsup x_n$.
    By construction of $y_k$, for each $k\in\n$, we may find $m_k\geq k$ such that
    \[
        x_{m_k} > y_k - \dfrac{1}{k}.
    \]
    Since $y_k \geq x_{m_k}$ and
    \[
        \lim_{k\to\infty}\left(y_k - \dfrac{1}{k}\right) = \lim y_k,
    \]
    the Squeeze Theorem implies that $x_{m_k}\to L^\upstar$.
    The sequence $(x_{m_k})$ is not necessarily a subsequence since $(m_k)$ may not be increasing.
    This is fixed using the lemma below.
\end{theorem}

\begin{lemma}{
Let $m_1, m_2,\ldots$ be a sequence of natural numbers such that $m_k\to\infty$.
Then there exists a subsequence $(m_{n_k})$ that is strictly increasing.
}
Let $n_1 \coloneqq 1$.
Suppose $n_1 < \cdots < n_j$ have been chosen.
Since $m_k\to\infty$, for any $C$ there exists $K$ such that $k \geq K \implies m_k > C$.
Take $C \coloneqq m_{n_j}$.
Define $n_{j+1} \coloneqq K$.
Then $m_{n_{j+1}} > m_{n_j}$.
\end{lemma}

\begin{example}
    \begin{enumerate}[label={(\arabic*)}]
        \item For any enumeration $(x_n)$ of $\q$, every real number is the limit of some subsequence of $(x_n)$.
        \item For any enumeration $(x_n)$ of $(0,1)$,
        \[
        \limsup x_n = 1,\quad \liminf x_n = 0.
        \]
    \end{enumerate}
\end{example}